\subsection{The isolated degree-twenty Bernoulli certificate}
\label{sec:bernoulli-residual-prime}

\begin{proposition}[K17-BERNOULLI-9496386083]
\label{prop:bernoulli-residual-prime}
\label{prop:bernoulli-residue}
For $p=9496386083$ and $n=p-11=9496386072$,
\[
                  B_n\equiv3753150554\ne0\pmod p.
 \tag{BC.1}
\]
This certifies the single class-group component required by
Lemma~\ref{lem:isolated-reflection}; it does not assert Vandiver's
conjecture for every character at this prime.
\end{proposition}

\begin{proof}
We give both the finite formula and its exact exhaustive verification.
The Bernoulli recurrence
\[
 \sum_{j=0}^{m}\binom{m+1}{j}B_j=0\qquad(m\ge1),\qquad B_0=1,
\]
shows inductively that $B_m$ is $p$-integral for $m\le p-2$,
because the coefficient $m+1$ of $B_m$ is a $p$-adic unit.
Faulhaber's polynomial identity then implies, for
$2\le n\le p-3$,
\[
                  T_n:=\sum_{a=1}^{p-1}a^n
                       \equiv pB_n\pmod{p^2}.
 \tag{BC.2}
\]
For $g$ prime to $p$, write $ga=r_a+pq_a$ with
$1\le r_a\le p-1$ and $q_a=\lfloor ga/p\rfloor$.
The $r_a$ permute the nonzero residues, so
\[
 T_n=\sum r_a^n
 \equiv g^nT_n-np g^{n-1}\sum q_a a^{n-1}\pmod{p^2}.
\]
If $g^n\ne1\pmod p$, division by $p$ and (BC.2) give
\[
 B_n\equiv
 \frac{n g^{n-1}}{g^n-1}
 \sum_{a=1}^{p-1}\left\lfloor\frac{ga}{p}\right\rfloor a^{n-1}
 \pmod p.
 \tag{BC.3}
\]
This derivation fixes the normalization and sign of the formula.

Both exact programs prove by trial division that $p$ and the factors in
\[
                     p-1=2\cdot17\cdot293\cdot953261
\]
are prime. For their respective choices $g=2$ and $g=5$, they check
$g^{(p-1)/q}\ne1\pmod p$ for every displayed prime divisor $q$.
Thus $g$ has order $p-1$. Put $h=(p-1)/2$ and let
$a_i$ be the least positive residue of $g^i$, for $0\le i<h$.
These powers give one member of each pair $\{a,-a\}$.
Since $n$ is even and
$q_{p-a}=g-1-q_a$, the sum in (BC.3) equals
\[
 S_g=\sum_{i=0}^{h-1}
       \left(2\left\lfloor\frac{ga_i}{p}\right\rfloor-(g-1)\right)w_i,
 \qquad w_i=a_i^{n-1}=a_i^{-11}.
 \tag{BC.4}
\]
The recurrence
$a_{i+1}=ga_i\bmod p$,
$w_{i+1}=g^{-11}w_i\bmod p$
computes it in constant memory with exactly
$h=4748193041$ terms.

For $g=2$, the signed coefficients in (BC.4) are $\pm1$,
$2^{-11}=2916614671\bmod p$, and the conversion factor simplifies to
\[
 \frac{n2^{n-1}}{2^n-1}
 =\frac{(-11)2^{-11}}{2^{-10}-1}
 =\frac{11}{2046}
 =5565086468\pmod p.
 \tag{BC.5}
\]
The first program evaluates the sum with exact Montgomery reduction.
For completeness, with $R=2^{64}$ and $v=-p^{-1}\bmod R$, its product
routine calculates
\[
 t=ab,\qquad m=(t\bmod R)v\bmod R,\qquad u=(t+mp)/R
\]
and subtracts $p$ if $u\ge p$.
The numerator is divisible by $R$,
$u\equiv abR^{-1}\pmod p$, and $0\le u<2p$.
All products use unsigned 128-bit integers and
\[
                 t+mp<p^2+2^{64}p<2^{100}<2^{128}.
\]
Unsigned 64-bit wrap is used deliberately only for reduction modulo
$2^{64}$ in the inverse and multiplier.
The orbit and signed sum updates are below $2p<2^{35}$.
The loop invariant is that the three state variables encode
$a_i$, $w_i$, and the first $i$ summands of (BC.4).
Thus the complete finite loop computes exactly the stated sum.
The actual full output is
\[
                         S_2=4849818985.
\]
The program also checks the unsimplified factor in (BC.3) against
(BC.5), and checks the half-orbit endpoint $a_h=w_h=-1$.

An independently written second program uses $g=5$, a different orbit,
and ordinary unsigned 128-bit modular division. Its signed coefficients
are $-4,-2,0,2,4$; the orbit update is below $5p<2^{36}$,
each coefficient times a weight is below $4p<2^{36}$, and the
weight product is below $p^2<2^{68}$. It independently verifies the
same primality and primitive-order conditions, performs all $h$ terms,
and checks $a_h=w_h=-1$. The completed outputs are
\[
\begin{array}{c|r|r|r}
g&S_g&n g^{n-1}/(g^n-1)\pmod p&B_n\pmod p\\ \hline
2&4849818985&5565086468&3753150554\\
5&7110767128&9130328994&3753150554.
\end{array}
\]
This proves (BC.1) by two exact exhaustive calculations of (BC.3).

The complete sources and reports are
\texttt{hecke\_general/bernoulli\_9496386083.cpp/.json} and
\texttt{hecke\_general/bernoulli\_root\_g5.cpp/.json}.
They use no numerical library and no floating-point arithmetic in
residues, exponents, loop counts, or comparisons.
Floating-point numbers occur only in timing and progress output.
The two completed runs took approximately $20.65$ and $41.09$ seconds;
neither is a sampled sum or a truncated benchmark.
\end{proof}
